\documentclass[11pt]{article}
\usepackage[margin=1in]{geometry}
\usepackage{amsmath,amssymb,amsthm}
\usepackage{lmodern}
\usepackage[colorlinks=true,linkcolor=blue,urlcolor=blue]{hyperref}
\newtheorem{theorem}{Theorem}
\newtheorem{lemma}[theorem]{Lemma}
\newcommand{\Z}{\mathbb Z}
\newcommand{\Q}{\mathbb Q}
\title{Weight preservation for modular Fibonacci recurrence periods\\
\large A formal verification of the orbit-counting argument}
\author{Myeongjin Kang}
\date{25 October 2026}
\setlength{\emergencystretch}{2em}
\begin{document}
\maketitle
\begin{abstract}
For positive integers $d\mid m$, reduction modulo $d$ preserves the total normalized
weight of every period of the Fibonacci recurrence and its parity transform.
We give the complete finite-orbit counting proof, including fixed points, and
formalize it in Lean. The argument is already present in a MathDB proof posted by
Shivam Patel; no novelty or priority is claimed. The formal theorem applies more
generally to $a_{n+2}=a_n+\varepsilon a_{n+1}$ for every integer $\varepsilon$.
\end{abstract}

\section{Statement and attribution}
Pudelko's Conjecture 17 in version 2 of \emph{Modular Periodicity of Random
Initialized Recurrences} \cite{pudelko} defines the weight of a period of length
$\ell$ modulo $s$ as $\ell/s^2$ and asserts conservation of weight under reduction
from a composite modulus to a divisor. Periods are considered up to cyclic shift,
as in the introduction of that source. In particular the constant zero sequence
is a period of length one. The theorem below includes every positive modulus,
so it implies the conjecture's composite-modulus case.

The orbit-fiber proof used here is due to the existing MathDB post by
\textbf{Shivam Patel}, dated 20 August 2026, solution
\nolinkurl{b64f2547-2755-408d-8ca3-997917a47b33} to problem 371201 \cite{patel}.
This note supplies a checked implementation of that argument, including its
identification of actual recurrence periods and their lifts. It makes no
new-to-literature claim. It does not address other conjectures in Pudelko's paper.

Fix $\varepsilon\in\Z$ and a positive integer $s$. On
$X_s=(\Z/s\Z)^2$ put
\[
T_s(a,b)=(b,a+\varepsilon b).
\]
Let $\mathcal P_s$ be the set of its permutation orbits, including singleton
orbits. Write $\ell(p)=|p|$ and $w_s(p)=\ell(p)/s^2$. For $d\mid m$ let
$\pi:X_m\to X_d$ be coordinatewise reduction.

\begin{theorem}\label{thm:main}
For every positive $m,d$ with $d\mid m$, every $\varepsilon\in\Z$, and every
$q\in\mathcal P_d$,
\[
 \sum_{\substack{p\in\mathcal P_m\\\pi(p)=q}}
       \frac{\ell(p)}{m^2}
   =\frac{\ell(q)}{d^2}.
\]
In particular this holds for the Fibonacci recurrence ($\varepsilon=1$) and
its parity transform ($\varepsilon=-1$).
\end{theorem}

\section{A finite equivariant counting lemma}
The group structure below is used only to establish uniform fibers.
The permutations need not preserve the group operations.

\begin{lemma}\label{lem:general}
Let $G,H$ be finite additive groups, let $f:G\to H$ be a surjective
homomorphism, and let $T$ and $S$ be permutations of $G$ and $H$ satisfying
$fT=Sf$. If $\mathcal P(T),\mathcal P(S)$ are their sets of orbits, then $f$
induces a map $\bar f:\mathcal P(T)\to\mathcal P(S)$ and, for every
$q\in\mathcal P(S)$,
\[
 \sum_{\substack{p\in\mathcal P(T)\\\bar f(p)=q}}
 \frac{|p|}{|G|}=\frac{|q|}{|H|}.
\]
\end{lemma}
\begin{proof}
Induction gives $f(T^n x)=S^n f(x)$ for every nonnegative integer $n$.
On a finite permutation orbit, every orbit element is a nonnegative iterate.
Consequently, $f$ maps each $T$-orbit into a single $S$-orbit, so $\bar f$ is
well defined. In fact its image is that entire orbit: if $y=S^n f(x)$ then
$T^n x$ lies in the orbit of $x$ and maps to $y$.

Let $k=|f^{-1}(0)|$. Surjectivity and translation in the group identify each
fiber $f^{-1}(y)$ with $f^{-1}(0)$; thus all fibers have cardinality $k$.
Counting the fibers over all of $H$ gives $|G|=|H|k$. Counting the fibers over
the subset $q$ gives $|f^{-1}(q)|=|q|k$.

The preimage $f^{-1}(q)$ is a disjoint union of exactly the $T$-orbits whose
image under $\bar f$ is $q$. To see both directions, a point in $f^{-1}(q)$
has its orbit mapped into $q$ by equivariance, whereas every point of an orbit
mapped to $q$ is in $f^{-1}(q)$. The quotient projection partitions this set
without omitting or duplicating any orbit. Hence
\[
 \sum_{\bar f(p)=q}|p|=|f^{-1}(q)|=|q|k.
\]
The groups are nonempty, so their orders are positive. Divide by
$|G|=|H|k$ to obtain the stated identity.
\end{proof}

\section{Application to the recurrence}
\begin{proof}[Proof of Theorem~\ref{thm:main}]
The map $T_s$ is a permutation for every $\varepsilon$, with inverse
\[
 T_s^{-1}(a,b)=(b-\varepsilon a,a).
\]
Both compositions are the identity by direct substitution. In particular no
coprimality hypothesis on $\varepsilon$ is required. A sequence satisfying
$a_{n+2}=a_n+\varepsilon a_{n+1}$ has state $(a_n,a_{n+1})=T_s^n(a_0,a_1)$.
A return of the state is equivalent to a period of the scalar sequence: a
period returns the consecutive pair, and a returned pair determines the
entire sequence by the recurrence and its inverse. Thus cyclic equivalence
of sequences is exactly membership in one $T_s$-orbit. The orbit's cardinality
is its least positive return time. This also covers fixed points, whose
cardinality and period are both one.

Coordinatewise reduction $\pi$ is a surjective homomorphism because $d\mid m$.
Reduction preserves addition, multiplication, and the integer coefficient, so
\[
 \pi\bigl(T_m(a,b)\bigr)
   =\bigl(\bar b,\bar a+\varepsilon\bar b\bigr)
   =T_d\bigl(\pi(a,b)\bigr).
\]
Lemma~\ref{lem:general} applies with $G=X_m$ and $H=X_d$. Their cardinalities
are $m^2$ and $d^2$, which gives the theorem. Equivalently, each fiber has
$(m/d)^2$ elements, so the total length of all orbits lifting $q$ is
$\ell(q)(m/d)^2$.
\end{proof}

\section{Formal implementation and audit}
The self-contained Lean source is \nolinkurl{formal/FibonacciWeight.lean}.
The period type is the quotient by
\nolinkurl{Equiv.Perm.SameCycle.setoid}, not a list of nontrivial cycles or a
permutation's support. The main theorem is
\nolinkurl{FibonacciWeight.fibonacci_weight_preservation}.
Its only inputs are $m,d$ with instances asserting they are nonzero, a proof
of $d\mid m$, an integer coefficient, and a downstairs quotient class.
In particular it does not assume the necessary orbit or fiber facts.

The following declarations provide the principal bridges:
\begin{itemize}
\item \nolinkurl{period_eq_iff_iterate} identifies quotient equality with time shift;
\nolinkurl{orbitSize_eq_minimalPeriod} identifies orbit cardinality with the least
positive return time.
\item \nolinkurl{periodMap} constructs the induced quotient map;
\nolinkurl{orbit_maps_onto} proves that every upstairs orbit maps onto its target.
\item \nolinkurl{orbit_sum_eq_preimage} partitions the actual preimage into
upstairs quotient classes. Its proof uses the finite fiberwise cardinal sum.
\item \nolinkurl{uniform_fibers}, \nolinkurl{preimage_card}, and \nolinkurl{total_card}
prove the counting identities from a surjective group homomorphism.
\item \nolinkurl{reduction_surjective} and \nolinkurl{reduction_commutes} prove the
hypotheses for the concrete maps; \nolinkurl{reduction_fiber_card} derives the
explicit square fiber cardinality.
\item \nolinkurl{zero_period_length} verifies that the zero fixed point has
length one.
\end{itemize}

The pinned compiler is Lean 4.34.1 and the mathlib revision is
\begin{center}\small\texttt{d13f23b723b8a846827a245b89c10fc7d3f11612}\end{center}
The local rebuild script checks the dependency revisions and their tracked
cleanliness, removes the local proof objects, compiles both the proof and
its declaration checks with warnings treated as errors, and audits all
printed dependencies against the allowlist
\nolinkurl{propext}, \nolinkurl{Classical.choice}, and \nolinkurl{Quot.sound}.
The release includes \nolinkurl{verify_outcome.py}; run it on the package directory with the pinned dependencies. It creates receipts under \nolinkurl{build/}.
No unchecked proof shortcut or additional axiom is used.

An independent review is included separately as \nolinkurl{review.md}. The formalization here
covers the quadratic integer-coefficient recurrence family above, not the
additional general matrix assertion in the posted proof.

\begin{thebibliography}{9}
\bibitem{pudelko}
Marc T. Pudelko.
\emph{Modular Periodicity of Random Initialized Recurrences}.
arXiv:2510.24882v2, 3 November 2025, Conjecture 17, equation (19).
\url{https://arxiv.org/html/2510.24882v2}.
\bibitem{patel}
Shivam Patel.
Proof posted to MathDB problem 371201, \emph{Weight preservation for Fibonacci
recurrences}, 20 August 2026.
Solution identifier \nolinkurl{b64f2547-2755-408d-8ca3-997917a47b33}.
Source record: \nolinkurl{source.json}.
\end{thebibliography}
\end{document}
